\documentclass[11pt,a4paper]{article}

\usepackage{amsmath, amssymb}
\usepackage{geometry}
\usepackage{listings}
\usepackage{xcolor}
\usepackage{hyperref}
\usepackage{booktabs}
\usepackage{array}
\usepackage{parskip}

\geometry{margin=2cm, top=2cm, bottom=2cm}

\setlength{\parskip}{3pt}
\setlength{\parindent}{0pt}

\hypersetup{
    colorlinks=true,
    linkcolor=blue,
    urlcolor=blue
}

\lstset{
    language=Python,
    basicstyle=\ttfamily\small,
    keywordstyle=\color{blue},
    commentstyle=\color{gray},
    stringstyle=\color{orange},
    showstringspaces=false,
    backgroundcolor=\color{gray!10},
    frame=single,
    framerule=0.4pt,
    breaklines=true
}

\title{\vspace{-1cm}Binary Image Classification Report:\\
K-Nearest Neighbors, Logistic Regression \& Gaussian Na\"{i}ve Bayes}
\author{}
\date{}

\begin{document}

\maketitle
\vspace{-0.8cm}
\tableofcontents
\newpage

%% ============================================================
\section{Introduction}
%% ============================================================

This report presents the implementation and evaluation of three machine learning algorithms --- \textbf{K-Nearest Neighbors (KNN)}, \textbf{Logistic Regression}, and \textbf{Gaussian Na\"{i}ve Bayes} --- for binary image classification on the MNIST handwritten digit dataset, as part of Phase~1 of the CSE382 Major Project. Each model is built from scratch and evaluated under a shared preprocessing pipeline using accuracy, precision, recall, F1-score, and confusion matrices.

%% ============================================================
\section{Data Preprocessing}
%% ============================================================

\subsection{Normalization}
Raw MNIST pixel values in $[0,255]$ are normalized to $[0,1]$: $\;\hat{x} = x / 255.0$

\subsection{Feature Extraction}
Three representations are compared: \textbf{Flatten} — each $28\times28$ image reshaped to $x \in \mathbb{R}^{784}$; \textbf{PCA} — projection onto the top $k$ eigenvectors retaining 95\% variance, fit on training data only; \textbf{HOG} — Histogram of Oriented Gradients with $4\times4$ cells, $2\times2$ blocks, and 9 orientation bins.

\subsection{Class Imbalance}
Random undersampling is applied to the training split only: $N_{\text{balanced}} = \min(|C_0|, |C_1|)$

\subsection{Feature Standardization}
Features are standardized using training-set statistics (zero-variance dimensions set to $\sigma=1$):
\[
z = \frac{x - \mu_{\text{train}}}{\sigma_{\text{train}}}
\]

\subsection{Train / Validation / Test Split}
The binary subset ($\sim$14,780 samples, seed = 42) is split as follows:

\vspace{4pt}
\begin{center}
\begin{tabular}{lcc}
\toprule
\textbf{Split} & \textbf{Proportion} & \textbf{Approx.\ Samples} \\
\midrule
Training   & 60\% & $\sim$8,868 \\
Validation & 20\% & $\sim$2,956 \\
Test       & 20\% & $\sim$2,956 \\
\bottomrule
\end{tabular}
\end{center}

%% ============================================================
\section{K-Nearest Neighbors (KNN)}
%% ============================================================

\subsection{Overview}
KNN is a non-parametric, instance-based classifier with no explicit training phase — the dataset is memorised and all computation occurs at prediction time. Given a query sample, KNN computes its distance to every training point, selects the $K$ closest neighbours, and assigns the majority class label. $K$ is tuned via 5-fold cross-validation over $K \in \{1, 3, 5, 7, 9, 11\}$.

\subsection{Mathematical Formulation}

\subsubsection*{1.\ Euclidean Distance (Vectorized)}
Given query matrix $X_q\,(\ell \times d)$ and training matrix $X_{\text{tr}}\,(m \times d)$, squared distances are computed via a single matrix multiplication:
\[
D^2 = \|X_q\|^2 + \|X_{\text{tr}}\|^2 - 2\,X_q \cdot X_{\text{tr}}^\top
\]

\subsubsection*{2.\ K-Selection via Argpartition}
The $K$ nearest indices are found in $O(n)$ rather than $O(n \log n)$:
\[
\text{k\_idx} = \operatorname{argpartition}(d_i,\;K)[:K]
\]

\subsubsection*{3.\ Majority Vote (Classification Rule)}
The predicted class is the most frequent label among the $K$ neighbours:
\[
\hat{y} = \operatorname*{arg\,max}_{c}\ \bigl|\{\,x_i \in \mathcal{N}_K(x) : y_i = c\,\}\bigr|
\]

\subsubsection*{4.\ Class Frequency Ratio (Probability Estimation)}
Soft probabilities are the fraction of the $K$ neighbours belonging to each class:
\[
\hat{P}(Y = c \mid x) = \frac{\text{count}_c}{K}
\]

\subsection{Results}

\vspace{4pt}
\begin{center}
\begin{tabular}{lcccc}
\toprule
\textbf{Feature} & \textbf{Best $K$} & \textbf{Val Acc.} & \textbf{Test Acc.} & \textbf{F1-Score} \\
\midrule
\textbf{Flatten} & \textbf{3} & \textbf{0.9994} & \textbf{0.9994} & \textbf{0.9994} \\
PCA (95\%)       & 5          & 0.9991          & 0.9991          & 0.9991          \\
HOG              & 5          & 0.9989          & 0.9989          & 0.9989          \\
\bottomrule
\end{tabular}
\end{center}
\vspace{4pt}

\noindent\textbf{Confusion Matrix — Flatten, $K=3$:}
\[
\begin{pmatrix} 1135 & 1 \\ 1 & 1172 \end{pmatrix}
\]

The best configuration is \textbf{Flatten with $K=3$} (99.94\% accuracy, F1 = 0.9994). Digits 0 and 1 are geometrically distinct, making raw pixel distances sufficient for near-perfect separation without dimensionality reduction.

\newpage

%% ============================================================
\section{Logistic Regression}
%% ============================================================

\subsection{Sigmoid Activation}

The predicted probability is computed via the sigmoid function:
\[
\hat{y} = \sigma(z) = \frac{1}{1 + e^{-z}} \qquad (B \times 1)
\]
Applied element-wise, $\hat{y}_i \in (0,1)$ is the predicted probability that sample $i$ belongs to digit $d$.

\subsection{Loss Function (Binary Cross-Entropy)}
For a mini-batch of size $B$:
\[
\mathcal{L}(w) = -\frac{1}{B} \sum_{i=1}^{B} \left[ y_i \log(\hat{y}_i) + (1 - y_i)\log(1 - \hat{y}_i) \right]
\]
The validation loss uses $N_{\text{val}} = 10{,}000$ samples and is computed after each full epoch to monitor generalisation.

\subsection{Handling Class Imbalance}
To account for an imbalanced class distribution in the training data, class-specific weights are introduced into the loss function. Given $N$ total samples, $n_{\text{pos}}$ positive examples, and $n_{\text{neg}}$ negative examples, the weights are:
\[
w_{\text{pos}} = \frac{N}{2\, n_{\text{pos}}}, \qquad w_{\text{neg}} = \frac{N}{2\, n_{\text{neg}}}
\]
These factors upweight the minority class and downweight the majority class, ensuring balanced gradient signals. The weighted binary cross-entropy loss becomes:
\[
\mathcal{L}(w) = -\frac{1}{B} \sum_{i=1}^{B} \left[ w_{\text{pos}}\, y_i \log(\hat{y}_i) + w_{\text{neg}}(1 - y_i)\log(1 - \hat{y}_i) \right]
\]

\subsection{Gradient Derivation via Chain Rule}
\textbf{Why} $\displaystyle\frac{\partial \mathcal{L}}{\partial w} = \frac{1}{B} X_b^\top (\hat{y} - y)$

We derive the gradient from first principles using the chain rule through three composed functions.

\noindent\textbf{Step 1 — Loss w.r.t.\ $\hat{y}_i$}
\[
\frac{\partial \mathcal{L}}{\partial \hat{y}_i}
= -\frac{1}{B}\left(\frac{y_i}{\hat{y}_i} - \frac{1-y_i}{1-\hat{y}_i}\right)
= -\frac{1}{B} \cdot \frac{y_i - \hat{y}_i}{\hat{y}_i(1-\hat{y}_i)}
\]

\noindent\textbf{Step 2 — $\hat{y}_i$ w.r.t.\ $z_i$ (sigmoid derivative)}
\[
\frac{\partial \hat{y}_i}{\partial z_i} = \sigma(z_i)\bigl(1 - \sigma(z_i)\bigr) = \hat{y}_i(1 - \hat{y}_i)
\]

\textit{Proof.} Let $\sigma = (1+e^{-z})^{-1}$:
\[
\frac{d\sigma}{dz}
= \frac{e^{-z}}{(1+e^{-z})^2}
= \frac{1}{1+e^{-z}} \cdot \frac{e^{-z}}{1+e^{-z}}
= \sigma \cdot (1 - \sigma)
\]

\noindent\textbf{Step 3 — $z_i$ w.r.t.\ $w$}
\[
\frac{\partial z_i}{\partial w} = x_i' \quad (785 \times 1)
\]

\noindent\textbf{Chaining Steps 1--3}
\[
\frac{\partial \mathcal{L}}{\partial w}
= \sum_{i=1}^{B}
\underbrace{\frac{\partial \mathcal{L}}{\partial \hat{y}_i}}_{\text{Step 1}}
\cdot
\underbrace{\frac{\partial \hat{y}_i}{\partial z_i}}_{\text{Step 2}}
\cdot
\underbrace{\frac{\partial z_i}{\partial w}}_{\text{Step 3}}
\]
Substituting:
\[
\frac{\partial \mathcal{L}}{\partial w}
= \sum_{i=1}^{B}
\left(-\frac{1}{B} \cdot \frac{y_i - \hat{y}_i}{\hat{y}_i(1-\hat{y}_i)}\right)
\cdot \hat{y}_i(1-\hat{y}_i) \cdot x_i'
\]
The $\hat{y}_i(1-\hat{y}_i)$ terms cancel exactly:
\[
\frac{\partial \mathcal{L}}{\partial w}
= \frac{1}{B} \sum_{i=1}^{B} (\hat{y}_i - y_i)\, x_i'
\]
Writing this as a matrix product:
\[
\frac{\partial \mathcal{L}}{\partial w}
= \frac{1}{B}\, X_b^\top (\hat{y} - y)
\qquad
\underbrace{(785 \times B)}_{X_b^\top} \cdot \underbrace{(B \times 1)}_{(\hat{y}-y)} = \underbrace{(785 \times 1)}_{\nabla w}
\]
This matches \texttt{de\_dw = np.dot(xi\_batch.T, error) / batch\_size} exactly, where \texttt{error = y\_hat - y}.

\subsection{Weight Update Rule}

\[
w \leftarrow w - \eta \cdot \frac{\partial \mathcal{L}}{\partial w} \qquad (785 \times 1)
\]
Matching \texttt{w -= learning\_rate * de\_dw}, where $\eta$ is \texttt{learning\_rate}.

\subsection{Mini-Batch Training}

Each epoch proceeds as follows:

\begin{enumerate}
    \item \textbf{Shuffle:} draw a random permutation of the 50,000 training indices ---
    \texttt{indices = np.random.permutation(N)}.

    \item \textbf{Batch loop:} for batch $k = 0, 1, \ldots, \lfloor N/B \rfloor - 1$:
    \[
    \text{start} = k \cdot B, \quad \text{end} = \text{start} + B
    \]
    Extract $X_b = X_{\text{shuffled}}[\text{start}:\text{end}]\ (B\times 785)$ and $y_b\ (B\times 1)$,
    then perform the forward pass and weight update described above.
   
\end{enumerate}

%% ============================================================
\section{Gaussian Na\"{i}ve Bayes}
%% ============================================================

\subsection{Methodology and Preprocessing}
The objective was to implement a Gaussian Na\"{i}ve Bayes classifier from scratch for binary classification. Digits 3 and 8 were selected due to their structural similarity.

\noindent\textbf{Preprocessing:}
\begin{itemize}
    \item \textbf{Filtering:} The raw MNIST dataset was filtered to include only instances of the target classes.
    \item \textbf{Flattening:} Each $28\times28$ pixel matrix was transformed into a 784-dimensional feature vector $x \in \mathbb{R}^{784}$.
    \item \textbf{Normalisation:} Pixel intensities were scaled from $[0, 255]$ to $[0, 1]$ to ensure numerical stability during probability density calculations.
\end{itemize}

\subsection{Mathematical Formulation}
The classifier predicts the class $y$ by maximising the posterior probability $P(y \mid x)$ via Bayes' Theorem:
\[
P(y \mid x) \propto P(y) \prod_{i} P(x_i \mid y)
\]
To prevent numerical underflow, the implementation uses log probabilities:
\[
\hat{y} = \operatorname*{arg\,max}_{c} \left[ \ln P(y=c) + \sum_{i} \ln P(x_i \mid y=c) \right]
\]
Each pixel's likelihood is modelled using the Gaussian distribution:
\[
P(x_i \mid y=c) = \frac{1}{\sqrt{2\pi\sigma^2}} \exp\!\left(-\frac{(x_i - \mu)^2}{2\sigma^2}\right)
\]

\noindent\textit{Implementation note:} A variance smoothing factor $\varepsilon = 10^{-2}$ was added to $\sigma^2$ to prevent division by zero.

\subsection{Experimental Results and Evaluation}

\begin{center}
\begin{tabular}{ll}
\toprule
\textbf{Metric} & \textbf{Value} \\
\midrule
Training Accuracy & 91.55\% \\
Testing Accuracy  & 92.59\% \\
\bottomrule
\end{tabular}
\end{center}

\noindent\textbf{Confusion Matrix:}
\[
\begin{pmatrix} 905 & 105 \\ 42 & 932 \end{pmatrix}
\]

\end{document}